% Source: book p. 181 (PDF 195).
\chapter{Inner Product Spaces}

In making the definition of a vector space, we generalized the linear structure
(addition and scalar multiplication) of $\R^2$ and $\R^3$. We ignored geometric
features such as the notions of length and angle. These ideas are embedded in
the concept of inner products, which we will investigate in this chapter.

Every inner product induces a norm, which you can think of as a length. This
norm satisfies key properties such as the Pythagorean theorem, the triangle
inequality, the parallelogram equality, and the Cauchy--Schwarz inequality.

The notion of perpendicular vectors in Euclidean geometry gets renamed to
orthogonal vectors in the context of an inner product space. We will see that
orthonormal bases are tremendously useful in inner product spaces. The
Gram--Schmidt procedure constructs such bases. This chapter will conclude by
putting together these tools to solve minimization problems.

\begin{definition*}[Standing assumptions for this chapter]\leavevmode
\begin{itemize}
\item $\F$ denotes $\R$ or $\C$.
\item $V$ and $W$ denote vector spaces over $\F$.
\end{itemize}
\end{definition*}
